% =====================================================================
%  Completed introduction.  Uses natbib (\citep / \citet).
%  New bibliography entries are in references_new.bib -- append them to
%  references.bib before compiling.
% =====================================================================

\section{Introduction}

The epoch of reionization (EoR) was the global phase transition of the
primordial neutral hydrogen into an ionized state. This process started with
the cosmic dawn (CD), when the first sources of stellar light were formed in
minihaloes of $\sim\!10^{5}$--$10^{6}\,M_{\odot}$ and the radiation emitted by
them began to couple, heat and locally ionize the surrounding medium, from
redshift $z \sim 30$ down to $z \sim 20$ \citep{Bromm2013,Klessen2023}. At the
other end, observations of the Ly$\alpha$ forest indicate that
reionization-related fluctuations were still present at $z \simeq 5.3$, and the
Thomson scattering optical depth of the CMB, $\tau = 0.054 \pm 0.007$,
constrains the midpoint of this process near $z \simeq 7.7$
\citep{Bosman2022,Planck2020}. Therefore, the range $5.5 \le z \le 20$ contains
the interval of cosmic time in which the first stellar sources altered the
thermodynamic state of the intergalactic medium (IGM)
\citep{Robertson2022,Klessen2023}.

What that interval does \emph{not} settle is which sources supplied the
ionizing photons, and whether they supplied them at the right time.
Star-forming galaxies are the standard answer \citep{Robertson2022}, but the
abundances and ionizing efficiencies inferred from \emph{JWST} are high enough
that, at face value and with canonical escape fractions, they would complete
reionization too early and overshoot both constraints quoted above --- the
so-called photon budget crisis \citep{Munoz2024}. A budget that is
over-supplied in photons and under-constrained in timing is precisely the
situation in which the secondary, non-stellar channels of energy injection
deserve to be quantified rather than assumed negligible.

Since then, the presence of cosmic rays (CR) is tightly linked to the
life-cycle of stellar populations. Stellar winds, accretion events, supernova
explosions, neutron stars and high-mass X-ray binaries (HMXBs) are
well-established accelerators of both hadronic and leptonic CRs: collisionless
shocks in the winds of massive stars and in their subsequent blast waves
develop the universal diffusive-shock-acceleration (DSA) power law, with
kinetic particle-in-cell simulations showing that electrons \emph{and} protons
are injected together at quasi-parallel shocks
\citep{Aharonian2019,Park2015}; the magnetized rotators left behind by the
explosions convert their spin-down power almost entirely into relativistic
pairs \citep{Bykov2017}; and accreting stellar-mass black holes drive outflows
whose X-ray output rises steeply towards low metallicity, so that the
metal-poor systems characteristic of the first galaxies are substantially more
luminous per unit star formation than local analogues
\citep{Mirabel2011,Kaur2022}. Because these engines all appear within a few
million years of the first star formation, CRs have been present throughout
the entirety of the EoR, interacting with the primordial IGM and injecting
energy into it \citep{Sazonov2015,Leite2017,Ruszkowski2023}.

For that injected energy to be a \emph{cosmological} agent, however, the
particles must first leave the object that accelerated them. Here the early
Universe is favourable: minihaloes have shallow potential wells, short magnetic
confinement times and low column densities, and the photoevaporation that
accompanies the first ionizing sources removes much of the confining gas, so
that CRs produced by the first supernova remnants are released essentially
directly into the surrounding neutral medium
\citep{Sazonov2015,Leite2017,Jana2018}. The relevant configuration is
therefore a relativistic particle propagating in a diffuse, cold and almost
fully neutral IGM, and the question becomes how it deposits its kinetic energy
there.

That deposition is not a single interaction but a cascade: high-energy CR
electrons and photons trigger extensive electromagnetic showers, in which each
generation of secondaries is itself energetic enough to ionize, excite and
produce further secondaries, so that the fate of the injected energy is decided
only after many generations \citep{Valdes2010,Slatyer2016}. The hierarchy of
the cascade is set by the energy of the particle. For high-energy electrons
the dominant energy loss mechanism is inverse Compton (IC) scattering off the
dense CMB, whose energy density grows as $(1+z)^{4}$ and therefore overwhelms
every other channel at these redshifts \citep{Blumenthal1970,Gould1972}. The
up-scattered photons, however, are only useful to reionization within a
restricted window: the parent electron must be energetic enough for the
scattered photon to exceed the $13.6\ \mathrm{eV}$ threshold, and it must not
be \emph{so} energetic that the resulting photon has a mean free path against
photoionization longer than the Hubble radius, in which case it free-streams
and its energy leaves the medium altogether \citep{Zdziarski1989,Slatyer2016}.
Locating these two boundaries quantitatively is one of the aims of this work.

For low-energy electrons the main cooling mechanism is instead the collisional
ionization of the cold neutral hydrogen, whose cross-section peaks a few times
above threshold and which converts the electron energy directly into free
electrons rather than into photons \citep{KimRudd1994,Kim2000}. Eventually the
thermalizing electron falls below $13.6\ \mathrm{eV}$ and can no longer ionize;
it then injects its remaining energy into the IGM, heating it via atomic
excitation and Coulomb collisions with the residual free electrons
\citep{Shull1985,Furlanetto2010}. Heating and ionization are thus not
independent processes but two branches of the same energy budget, fixed by the
same set of cross-sections. That the heating branch was indeed active is now
an observational statement: the HERA Phase~I 21-cm power-spectrum limits
require the IGM to have been warmed above the adiabatic cooling floor no later
than $z = 10.4$, excluding the whole family of ``cold reionization'' histories
\citep{HERA2023,Mertens2025}. What remains to be established is how much of the
same budget was diverted into the ionization branch.

The question is not new. \citet{Nath1993} asked directly whether CRs from young
galaxies could reionize the IGM, and it has been revisited with progressively
more complete microphysics ever since
\citep{Samui2005,Sazonov2015,Leite2017,Jana2018,Bera2023,GesseyJones2023}. The
literature is sharply asymmetric between the two branches: the heating is found
to be real and substantial, reaching $\Delta T \simeq 10$--$200\ \mathrm{K}$ by
$z = 10$ \citep{Sazonov2015,Leite2017}, whereas the ionization is concluded to
be negligible \citep{Leite2017}. That conclusion, however, has almost
invariably been reached for the \emph{hadronic} component, which carries the
bulk of the CR energy budget. The electronic component interacts through a
qualitatively different set of channels --- IC scattering on the CMB above all,
which at fixed total energy is suppressed for protons by
$(m_e/m_p)^{4} \simeq 10^{-13}$ and is therefore irrelevant for them, but which
dominates every other loss for electrons above $\sim\!100\ \mathrm{keV}$.
Whether the negligible-ionization verdict carries over to the leptonic
component is not obvious \emph{a priori}.

In this work, we extensively analyze and discuss the most relevant physical
effects and phenomena that affect the CR electron propagation and energy
deposition through the EoR. We focus mainly on the cooling mechanisms by which
CR electrons can heat and ionize the early IGM, evaluate their efficiency, and
determine the optimal parameters for which those electrons could have
contributed to the global reionization of the Universe. Finally, we discuss
which astrophysical scenarios could feasibly produce CR electrons to change the
thermodynamic state of the IGM during the EoR.
